\documentclass[12pt,letterpaper]{article}

\usepackage[margin=1in]{geometry}
\usepackage{amsmath,amssymb}
\usepackage{fancyhdr}

\setlength{\parindent}{0pt}
\setlength{\parskip}{0.8em}
\setlength{\headheight}{15pt}

\pagestyle{fancy}
\fancyhf{}
\fancyhead[L]{FIT3155 SAMPLE EXAM PAPER}
\fancyfoot[C]{Page \thepage{} of 22}
\renewcommand{\headrulewidth}{0pt}
\renewcommand{\footrulewidth}{0pt}

\newcommand{\marksbox}[1]{%
  \vfill
  \begin{flushright}
    \fbox{\parbox[c][1.1cm][c]{1.3cm}{\centering #1}}
  \end{flushright}
}

\newcommand{\questiontitle}[2]{%
  \textbf{Question #1}\hfill\textbf{(#2 marks)}
}

\newcommand{\workingpage}{%
  \newpage
  \begin{center}
    (Page intentionally left blank for working space)
  \end{center}
}

\begin{document}

\begin{titlepage}
  \thispagestyle{empty}
  \vspace*{0.18\textheight}
  {\Large FIT3155 \hfill SAMPLE EXAM PAPER\par}
  \vspace{2.2cm}
  \begin{center}
    {\LARGE\bfseries FIT3155 SAMPLE EXAM PAPER\par}
    \vspace{0.8cm}
    {\large Generated Sample Exam 01\par}
  \end{center}
\end{titlepage}

\setcounter{page}{2}

\section*{INSTRUCTIONS}

\begin{itemize}
  \item This exam is worth 60 marks.
\end{itemize}

\begin{center}
  \textbf{General exam technique}
\end{center}

\begin{itemize}
  \item Some students lose marks by not attempting all questions.

  \item Suppose you are likely to get 3/6 on a question after spending say 10 minutes on it.
  Spending another 10 minutes on the same question may gain you at most 3 additional
  marks. In contrast, spending those 10 minutes on a new question could earn you another
  6 marks. The key point is: do not get stuck on one question.

  \item A good strategy is to first answer the questions you are most confident about before
  attempting the others. In other words, answering questions strictly serially during an
  exam can be suboptimal.

  \item Do not waste time answering an unrelated question in the hope of gaining partial marks.
  Marking will be based on the question that was actually asked.

  \item Where necessary, justify your answers with clear and concise explanations, without being
  overly wordy.
\end{itemize}

\newpage

\questiontitle{1}{6}

Two strings $S[1 \ldots n]$ and $T[1 \ldots n]$ have the same length. We say that
$S$ is a cyclic rotation of $T$ if $S$ can be obtained by moving some prefix of
$T$ to the end of $T$.

Write pseudocode for an efficient algorithm that determines whether $S$ is a
cyclic rotation of $T$. If it is, your algorithm should output all starting
positions in $T$ that generate $S$ as a rotation. If it is not, your algorithm
should output ``not a rotation''.

For this problem, you may assume that \texttt{zalg} is available as a built-in
function that takes a string as input and returns the array of its Z-values. Do
not write pseudocode for \texttt{zalg}. State the worst-case time and space
complexity of your algorithm.

\marksbox{6}

\workingpage
\newpage

\questiontitle{2}{6}

In the Boyer--Moore algorithm for exact matching of the pattern
$pat[1 \ldots m]$ against the text $txt[1 \ldots n]$, consider a general
iteration in which $pat$ is aligned with $txt[j \ldots j+m-1]$.

Suppose a right-to-left scan results in a mismatch at position
$1 \leq k \leq m$, and suppose the good-suffix array value satisfies
$gs[k+1] = 0$. In this case the matched-prefix array is used, giving a shift of
$m - mp[k]$ positions to the right.

Prove that this matched-prefix shift is safe, i.e., such a shift can never skip
$pat$ past any valid occurrence in $txt$.

\marksbox{6}

\workingpage
\newpage

\questiontitle{3}{6}

Let $S[1 \ldots n-1]\$$ be a string whose Burrows--Wheeler Transform is
$L[1 \ldots n]$. Assume that the alphabet of $S$ is $\Sigma$, and that $\$$ is a
unique terminal symbol.

Assume the following two preprocessed structures are available:
\[
  rank(c) = \text{the number of characters in } S \text{ that are strictly
  smaller than } c,
\]
and
\[
  Occ(c,i) = \text{the number of occurrences of } c \text{ in } L[1 \ldots i].
\]

Write pseudocode for the backward-search algorithm that takes $L$ and a pattern
$pat[1 \ldots m]$ as input, and returns the number of occurrences of $pat$ in
$S$. State the worst-case time complexity of the search, assuming
$rank$ and $Occ$ queries take $O(1)$ time.

\marksbox{6}

\workingpage
\newpage

\questiontitle{4}{6}

These questions relate to Ukkonen's linear-time suffix tree construction
algorithm for a string $S[1 \ldots n]$.

\begin{enumerate}
  \item[(i)] In phase $i+1$, suppose suffix extension $j$ is completed by Rule 3.
  Explain why the remaining suffix extensions in the same phase do not need to
  be performed explicitly.

  \item[(ii)] Suppose suffix extension $j$ creates a new internal node whose
  path-label is $\alpha = c\beta$ using Rule 2 regular. Explain when and why a
  suffix link should be created from this new internal node to the internal node
  whose path-label is $\beta$.
\end{enumerate}

\marksbox{6}

\workingpage
\newpage

\questiontitle{5}{6}

\begin{enumerate}
  \item[(i)] Work out the Elias-Omega binary codeword for the positive integer
  $31$.

  \item[(ii)] The following tuples were derived using the LZ77 method of
  encoding:
  \[
    [0,0,a]\quad
    [0,0,b]\quad
    [2,2,c]\quad
    [5,4,d]\quad
    [3,3,e]
  \]

  Decode these LZ77 tuples to regenerate the original string.
\end{enumerate}

\marksbox{6}

\workingpage
\newpage

\questiontitle{6}{6}

Perform one iteration of the Miller--Rabin primality test on $n = 77$ using
base $a = 10$.

Your answer should:
\begin{itemize}
  \item express $n-1$ in the form $2^s t$, where $t$ is odd;
  \item compute the values checked during this Miller--Rabin iteration using
  modular exponentiation; and
  \item state whether this base proves that $77$ is composite, or whether this
  iteration is inconclusive.
\end{itemize}

\marksbox{6}

\workingpage
\newpage

\questiontitle{7}{6}

Write pseudocode for the \textsc{Decrease-Key} operation on a Fibonacci heap.
Your pseudocode should include any required \textsc{Cut} and
\textsc{Cascading-Cut} operations.

Then state the worst-case actual time and the amortized time of
\textsc{Decrease-Key}, using the potential function
\[
  \Phi(H) = \operatorname{nTrees}(H) + 2\operatorname{nMarked}(H).
\]

\marksbox{6}

\workingpage
\newpage

\questiontitle{8}{6}

A dynamic array starts empty. The first append creates capacity $1$. Whenever an
append is attempted on a full array, a new array with twice the previous capacity
is allocated, all existing elements are copied into it, and then the new element
is appended.

Using the aggregate method of amortized analysis, determine the amortized time
complexity of an append operation over a sequence of $n$ append operations.
Your answer should clearly account for both ordinary writes and copying work
caused by resizing.

\marksbox{6}

\workingpage
\newpage

\questiontitle{9}{6}

For any B-tree with minimum degree $t$ that contains $N$ keys, derive a lower
bound on the height of the B-tree in terms of $N$ and $t$.

State clearly which convention you are using for the height of the tree.

\marksbox{6}

\workingpage
\newpage

\questiontitle{10}{6}

Solve the following optimization problem using the tableau simplex algorithm.
If the optimal solution is not unique, state this clearly.

\[
\begin{array}{rl}
\text{Maximize} & z = x + 2y \\
\text{Subject to} & 2x + 4y \leq 12, \\
                  & 4x + 3y \leq 16, \\
                  & x \geq 0,\quad y \geq 0.
\end{array}
\]

\marksbox{6}

\workingpage

\vfill
\begin{center}
  End of Exam
\end{center}

\end{document}
